\subsection{施瓦茨函数}

我们用 \(\mathcal{S}(\mathbf{R}^{n})\) 表示无穷次可微函数 \(\varphi\) （在 \(\mathbf{R}^{n}\) 上、取复值）所成的空间，其中对所有 \(\alpha\) 与 \(\beta\) （属于 \(\mathbf{N}^{n}\) ），函数 \(\mathrm{X}^{\beta}\partial^{\alpha}\varphi\) 在 \(\mathbf{R}^{n}\) 上有界。我们给 \(\mathcal{S}(\mathbf{R}^{n})\) 赋以由可数半范数族 \((q_{\alpha,\beta})_{(\alpha,\beta)\in\mathbf{N}^{n}\times\mathbf{N}^{n}}\) 定义的局部凸拓扑，其中 \(q_{\alpha,\beta}\) 定义为

\[
q _ {\alpha , \beta} (\varphi) = \sup _ {x \in \mathbf {R} ^ {n}} | x ^ {\beta} \partial^ {\alpha} \varphi (x) | = \| \mathrm{X} ^ {\beta} \partial^ {\alpha} \varphi \| _ {\infty}
\]

（对 \(\varphi\in\mathcal{S}(\mathbf{R}^{n})\) ）。该拓扑是分离的。它也由半范数 \(\widetilde{q}_{\alpha,k}\) 定义，其中

\[
\widetilde {q} _ {\alpha , k} (\varphi) = \sup _ {x \in \mathbf {R} ^ {n}} \| x \| ^ {k} | (\partial^ {\alpha} \varphi) (x) |.
\]

对所有 \(\varphi\in\mathcal{S}(\mathbf{R}^{n})\) ，其中 \(k\in\mathbf{N}\) ，\(\alpha\in\mathbf{N}^{n}\) 。

我们称 \(\mathcal{S}(\mathbf{R}^{n})\) 为 \(\mathbf{R}^{n}\) 上的施瓦茨空间或施瓦茨函数空间。

注. --- 设 \(\varphi\in\mathcal{S}(\mathbf{R}^{n})\) 。对每个 \(k\in\mathbf{N}\) ，我们有

\[
\lim _ {\| x \| \to + \infty} \| x \| ^ {k} \varphi (x) = 0,
\]

因为函数 \(x \mapsto \|x\|^{k+1} \varphi(x)\) 有界。

例. --- 函数 \(\gamma_{n}\) ——在 \(\mathbf{R}^{n}\) 上由 \(\gamma_{n}(x)=\exp(-\|x\|^{2})\) 定义——属于 \(\mathcal{S}(\mathbf{R}^{n})\) 。事实上，对 k 作归纳可以证明：对每个整数 \(k\in \mathbf{N}\) ，存在多项式 \(P_{k}\in\mathbf{R}[X]\) 使得 \(\partial_{k}\gamma_{1}=P_{k}\gamma_{1}\) 。

于是，对所有 \(\alpha = (\alpha_{i})\in \mathbf{N}^{n}\) 与 \(\beta = (\beta_{i})\in \mathbf{N}^{n}\) ，以及每个 \(x = (x_{i})\in \mathbf{R}^{n}\) ，可得

\[
| (\mathrm{X} ^ {\beta} \partial^ {\alpha} \gamma_ {n}) (x) | = \prod_ {i = 1} ^ {n} | x _ {i} | ^ {\beta_ {i}} | \mathrm{P} _ {\alpha_ {i}} (x _ {i}) | \gamma_ {1} (x _ {i}),
\]

它在 x 于 \(\mathbf{R}^{n}\) 中变化时是有界量。

设 \(\alpha \in \mathbf{N}^n\) 与 \(\beta \in \mathbf{N}^n\) 。若 \(\varphi \in \mathcal{S}(\mathbf{R}^n)\) ，则 \(\mathrm{X}^{\beta}\partial^{\alpha}(\varphi)\) 仍是施瓦茨函数；如此定义的映射 \(\varphi \mapsto \mathrm{X}^{\beta}\partial^{\alpha}\varphi\) ——从 \(\mathcal{S}(\mathbf{R}^n)\) 到自身——是连续的。

空间 \(\mathcal{S}(\mathbf{R}^{n})\) 是一个拓扑代数。更确切地说，若 \(\varphi_{1}\) 与 \(\varphi_{2}\) 属于 \(\mathcal{S}(\mathbf{R}^{n})\) ，则 \(\varphi_{1}\varphi_{2}\) 是满足

\[
q _ {\alpha , \beta} (\varphi_ {1} \varphi_ {2}) \leqslant \sum_ {0 \leqslant \gamma \leqslant \alpha} \binom{\alpha}{\gamma} q _ {\gamma , \beta} (\varphi_ {1}) q _ {\alpha - \gamma , 0} (\varphi_ {2})\tag{2}
\]

（对所有 \(\alpha\) 与 \(\beta \in \mathbf{N}^n\) ；IV, p. 196，命题 1）的施瓦茨函数。

\(\mathcal{S}(\mathbf{R}^{n})\) 到空间 \(\mathcal{C}^{\infty}(\mathbf{R}^{n})\) （赋以 IV, p. 200 的 §3 中所述的拓扑）中的典范包含是连续的，因为

\[
\sup _ {x \in K} | \partial^ {\alpha} \varphi (x) | \leqslant q _ {\alpha , 0} (\varphi)
\]

对 \(\mathbf{R}^{n}\) 的每个紧子集 K、所有 \(\alpha \in \mathbf{N}^{n}\) 及每个施瓦茨函数 \(\varphi\) 都成立。

引理 7. --- 设 \(k \in \mathbf{N}\) ，\(\alpha \in \mathbf{N}^n\) 。对每个函数 \(\varphi \in \mathcal{S}(\mathbf{R}^n)\) 及每个实数 \(T > 0\)，我们有

\[
\widetilde {q} _ {\alpha , k} (\varphi) \leqslant \mathrm{T} ^ {k} \sup _ {\| x \| \leqslant \mathrm{T}} | \partial^ {\alpha} \varphi (x) | + \frac {1}{\mathrm{T}} \widetilde {q} _ {\alpha , k + 1} (\varphi).\tag{3}
\]

事实上，有

\[
\begin{array}{c} \widetilde {q} _ {\alpha , k} (\varphi) \leqslant \sup _ {\| x \| \leqslant T} \| x \| ^ {k} | \partial^ {\alpha} \varphi (x) | + \sup _ {\| x \| > T} \| x \| ^ {k} | \partial^ {\alpha} \varphi (x) | \\ \leqslant T ^ {k} \sup _ {\| x \| \leqslant T} | \partial^ {\alpha} \varphi (x) | + \frac {1}{T} \sup _ {\| x \| > T} \| x \| ^ {k + 1} | \partial^ {\alpha} \varphi (x) |. \end{array}
\]

命题 10. --- 设 B 是 \(\mathcal{S}(\mathbf{R}^{n})\) 的有界子集。\(\mathcal{S}(\mathbf{R}^{n})\) 在 B 上诱导的拓扑与 \(\mathcal{C}^{\infty}(\mathbf{R}^{n})\) 诱导的拓扑一致。

由于 \(\mathcal{S}(\mathbf{R}^{n})\) 到 \(\mathcal{C}^{\infty}(\mathbf{R}^{n})\) 中的包含是连续的，只需证明：对 \(\mathcal{S}(\mathbf{R}^{n})\) 的每个开子集 V，交集 \(V \cap B\) 对 \(\mathcal{C}^{\infty}(\mathbf{R}^{n})\) 诱导的拓扑而言在 B 中是开的。

设 V 是 \(\mathcal{S}(\mathbf{R}^{n})\) 的开子集。设 \(\varphi_{0} \in \mathrm{V} \cap \mathrm{B}\) 。存在有限集 I、族 \((\alpha_{i}, k_{i})_{i \in \mathrm{I}} \in (\mathbf{N}^{n} \times \mathbf{N})^{\mathrm{I}}\) 及实数 \(\varepsilon > 0\) ，使得 V 包含由 \(\varphi \in \mathcal{S}(\mathbf{R}^{n})\) ——满足

\[
\sup _ {i \in I} \widetilde {q} _ {\alpha_ {i}, k _ {i}} (\varphi - \varphi_ {0}) \leqslant \varepsilon .
\]

——所成之集。由于 B 在 \(\mathcal{S}(\mathbf{R}^{n})\) 中有界，存在 \(M > 0\) ，使得半范数 \(\widetilde{q}_{\alpha_{i},k_{i} + 1}\) （对 \(i\in I\) ）以 \(M\) 为界（在 \(B\) 上）。设 \(\delta >0\) 与 \(T > 0\) 为实数。依据不等式 (3)，只要 \(\varphi \in B\) 满足界

\[
\sup _ {i \in I} \sup _ {\| x \| \leqslant T} | \partial^ {\alpha_ {i}} (\varphi - \varphi_ {0}) | \leqslant \delta ,\tag{4}
\]

我们就有

\[
\sup _ {i \in I} \widetilde {q} _ {\alpha_ {i}, k _ {i}} (\varphi - \varphi_ {0}) \leqslant \delta T ^ {k} + \frac {2 M}{T}.
\]

取 \(T = \frac{4M}{\varepsilon}\) ，再取 \(\delta = \frac{\varepsilon}{2T^{k}}\) 。可以看出，\(V \cap B\) 包含 \(\varphi_{0}\) 在 B 中的、对 \(\mathcal{C}^{\infty}(\mathbf{R}^{n})\) 诱导拓扑而言由 (4) 定义的邻域。证明至此完成。

推论. --- 设 \((\varphi_{m})_{m\in\mathbf{N}}\) 是 \(\mathcal{S}(\mathbf{R}^{n})\) 中的有界序列。对每个函数 \(\varphi\in\mathcal{S}(\mathbf{R}^{n})\) ，下列断言等价：

\begin{enumerate}
\def\labelenumi{\alph{enumi})}
\item
  序列 \((\varphi_{m})\) 收敛于 \(\varphi\) （在 \(\mathcal{S}(\mathbf{R}^n)\) 中）；
\item
  序列 \((\varphi_{m})\) 收敛于 \(\varphi\) （在 \(\mathcal{C}^{\infty}(\mathbf{R}^{n})\) 中）。
\end{enumerate}

注. --- 序列 \((\varphi_{m})\) （在 \(\mathcal{C}^{\infty}(\mathbf{R}^{n})\) 中）收敛当且仅当：对每个 \(\alpha\in \mathbf{N}^{n}\) ，序列 \((\partial^{\alpha}\varphi_{m})\) 收敛于某个函数 \(\varphi^{(\alpha)}\) （在赋以紧收敛拓扑的 \(\mathcal{C}(\mathbf{R}^{n})\) 中）。这时我们有 \(\varphi^{(\alpha)}=\partial^{\alpha}\varphi\) ，且 \((\varphi_{m})\) 收敛于 \(\varphi^{(0)}\) 。

事实上，该条件是必要的。反之，若诸序列 \((\partial^{\alpha}\varphi_{m})\) 收敛于函数 \(\varphi^{(\alpha)}\) ——对所有 \(\alpha\in \mathbf{N}^{n}\) ——，则由 FVR, II, p. 2，定理 1 可得 \(\varphi^{(\alpha)}=\partial^{\alpha}\varphi^{(0)}\) ，这意味着序列 \((\varphi_{m})\) 收敛于 \(\varphi^{(0)}\) （在 \(\mathcal{C}^{\infty}(\mathbf{R}^{n})\) 中）。

命题 11. --- 空间 \(\mathcal{S}(\mathbf{R}^{n})\) 是弗雷歇空间和蒙泰尔空间。

由于空间 \(\mathcal{C}^{\infty}(\mathbf{R}^{n})\) 是完备的（EVT, III, p. 9，例 \(b\) ），命题 10 的推论表明 \(\mathcal{S}(\mathbf{R}^{n})\) 中的每个柯西序列都在 \(\mathcal{S}(\mathbf{R}^{n})\) 中收敛，因为它有界且在 \(\mathcal{C}^{\infty}(\mathbf{R}^{n})\) 中收敛。

于是空间 \(\mathcal{S}(\mathbf{R}^n)\) 是弗雷歇空间；特别地，它是桶型的（EVT, III, p. 25，命题 2 的推论）。设 B 是 \(\mathcal{S}(\mathbf{R}^n)\) 的有界子集，\((\varphi_m)_{m \in \mathbf{N}}\) 是取值于 B 的序列。由于 \(\mathcal{C}^\infty(\mathbf{R}^n)\) 是蒙泰尔空间（EVT, IV, p. 18，例 (4)），存在 \((\varphi_m)_{m \in \mathbf{N}}\) 的子序列在 \(\mathcal{C}^\infty(\mathbf{R}^n)\) 中收敛，从而在 \(\mathcal{S}(\mathbf{R}^n)\) 中收敛（命题 10）。因而 B 在 \(\mathcal{S}(\mathbf{R}^n)\) 中相对紧。由此得出 \(\mathcal{S}(\mathbf{R}^n)\) 是蒙泰尔空间。

